\documentclass[11pt,a4paper]{article}
\usepackage[T1]{fontenc}
\usepackage[utf8]{inputenc}
\usepackage[french]{babel}
\usepackage{lmodern}
\usepackage{amsmath,amssymb}
\usepackage[margin=2.6cm]{geometry}
\usepackage{booktabs}
\usepackage{enumitem}
\usepackage{hyperref}
\hypersetup{colorlinks=true,linkcolor=blue!50!black}

\newcommand{\NW}{\mathrm{NW}}
\newcommand{\code}[1]{\texttt{#1}}

\title{Spécification technique de l'implémentation M2\\
\large (workspace \code{m2\_fable})}
\author{Agent de réplication (workspace \code{m2\_fable})}
\date{3 juillet 2026}

\begin{document}
\maketitle

\begin{abstract}
Cette spécification fige, \emph{avant} l'implémentation, les structures de
données, les conventions laissées ambiguës par le document de conception, les
invariants, les tests, les formats de sortie et le protocole de validation.
Toute convention tranchée ici est marquée \textbf{[C$n$]} et sera rappelée
dans les rapports de résultats : ce sont des choix, pas des faits.
\end{abstract}

\section{Arborescence et environnement}

\begin{itemize}[nosep]
\item Code : \code{m2\_fable/src/m2/\{config, entities, contracts, market,
  bankruptcy, simulation, metrics, analysis, plots\}.py}
\item Expériences : \code{m2\_fable/experiments/m2/\{run\_baseline,
  run\_validation, run\_grid, run\_ablation\}.py}
\item Tests : \code{m2\_fable/tests/test\_*.py} (style pytest) + runner
  stdlib \code{tests/run\_all.py} (pytest absent du venv, installation
  interdite).
\item Python : \code{/home/anatole/jupyter/.venv/bin/python3}
  (numpy 2.4.3, scipy 1.17.1, matplotlib 3.10.8).
\item Résultats : \code{m2\_fable/experiments/m2/results/<run\_id>/}.
\end{itemize}

\section{Structures de données}

\paragraph{Entités (structure de tableaux).} Les entités sont stockées en
« struct of arrays » (listes Python homogènes, index $=$ identifiant entier
croissant, jamais réutilisé) :
\code{w[i]}, \code{alive[i]}, \code{birth[i]}, \code{claims[i]},
\code{debts[i]} (agrégats de contrats, maintenus incrémentalement),
\code{int\_in[i]}, \code{int\_out[i]} (flux d'intérêts du pas courant),
\code{defaulted[i]} (défaut de service au pas courant).
Justification : $10^3$--$10^4$ entités, accès par index dominant ; les
dataclasses par entité coûteraient cher sans rien apporter ici. Le module
\code{entities.py} encapsule ces tableaux dans une classe \code{Population}
avec méthodes \code{born()}, \code{kill()}, \code{nw(i)}.

\paragraph{Contrats.} \code{contracts.py} : classe \code{LoanBook}.
Un contrat $=$ enregistrement \code{[lender, borrower, q, r]} dans un dict
\code{loans[loan\_id]}, avec index inverses \code{by\_lender[i]} et
\code{by\_borrower[i]} (sets d'ids). Toutes les mutations passent par
\code{add(l,b,q,r)}, \code{remove(lid)}, \code{set\_principal(lid,q)} qui
maintiennent \code{claims}/\code{debts}. Aucun autre code ne touche à ces
agrégats.

\paragraph{Config.} \code{config.py} : dataclass gelée \code{M2Config}
(\code{alpha=1.0, delta=0.05, lam=10.0, k=6, sigma=0.25, w0=30.0, d0=28.0,
seed, T}, tolérances \code{tol\_nw=1e-9, q\_min=1e-9}, variantes nommées :
\code{rate\_rule='geom'|'arith'}, \code{fail\_lender\_loans='transfer'|'cancel'},
\code{fail\_residual='prorata'|'destroy'}, \code{credit=True|False}).
Sérialisée en JSON dans chaque run. Toute variante expérimentale passe par un
champ nommé de la config — jamais par une constante modifiée à la main.

\section{Conventions tranchées}

\begin{description}[style=nextline]
\item[C1 — Choc « commun en loi », i.i.d. par entité.] $\eta_i$ est tiré
  indépendamment pour chaque entité à chaque pas, de la même loi
  $\mathcal N(-\sigma^2/2,\sigma^2)$. (Lecture confirmée par le prototype ;
  un choc agrégé serait une variante distincte, non implémentée ici.)
\item[C2 — Ordre de service des intérêts.] Contrats parcourus par
  \code{loan\_id} croissant (ordre de création, déterministe). Paiement
  intégral si $w_b \ge rq$ ; sinon transfert de tout $w_b$, $w_b \gets 0$,
  marquage \code{defaulted[b]}. Les contrats suivants du même emprunteur
  reçoivent 0 (et le marquent aussi en défaut, idempotent). Alternative au
  prorata documentée, non retenue (coût : deux passes ; bénéfice attendu
  nul au premier ordre — à re-examiner si les défauts partiels s'avèrent
  fréquents).
\item[C3 — Exclusion des défaillants du marché.] Une entité marquée
  \code{defaulted} au pas courant ne participe pas au marché du crédit de ce
  pas (elle fera faillite en phase 7). Suit le prototype.
\item[C4 — Garde de gain à l'échange.] Un round ne crée un contrat que si
  $w_\ell > w_b$ strictement et $q \ge $ \code{q\_min}. Si $w_b = 0$,
  $r^*_b$ est évalué avec le garde $\max(w_b, 10^{-9})$.
\item[C5 — Faillite : créances de la faillie transférées au prorata.]
  Chaque contrat $(i, b, q, r)$ où $i$ est prêteuse est scindé entre les
  créanciers contractuels $c$ de $i$ au prorata de leurs créances $q_c$ :
  principal $q \cdot q_c / \sum q_c$, même taux, même emprunteur. Parts
  $<$ \code{q\_min} abandonnées (perte sèche de l'emprunteur… en fait
  réduction de sa dette : voir invariant I7). Si un créancier $c$ est
  l'emprunteur $b$ du contrat transféré, la part est annulée (dette envers
  soi-même). Variante \code{cancel} : contrats prêtés simplement annulés
  (c'est ce que fait le prototype E7c — ablation obligatoire, §7).
\item[C6 — Faillite : capital résiduel.] $w_i$ résiduel réparti aux
  créanciers contractuels \emph{vivants} au prorata ; s'il n'y a aucun
  créancier contractuel (seule $d_0$ au passif), le résiduel est détruit.
  Variante \code{destroy} : détruit dans tous les cas.
\item[C7 — Ordre de résolution des cascades.] File des insolvables triée par
  identifiant croissant ; on traite une faillite entièrement (annulations,
  transferts, distribution), puis on re-balaie les vivantes pour
  $\NW < -$\code{tol\_nw}. Itération jusqu'à point fixe. Les créanciers déjà
  morts au moment d'une distribution ne reçoivent rien (leur part est
  détruite).
\item[C8 — Revenu.] $y_i$ du pas $=$ montant d'extraction effectivement
  ajouté en phase 3 ($\alpha\sqrt{w_i}$ évalué après choc) $+$ intérêts
  reçus en phase 4. Mesuré dans le snapshot de fin de pas. Revenu net $=$
  $y_i -$ intérêts payés.
\item[C9 — Les nouveau-nés du pas participent à tout le pas] (choc,
  extraction, marché). Âge $=$ $t_{\mathrm{snapshot}} - t_{\mathrm{naissance}}$.
\item[C10 — RNG.] Un unique \code{numpy.random.default\_rng(seed)} par
  simulation ; aucun usage de \code{random} global. Ordre des tirages fixé
  par la séquence des phases (reproductibilité bit à bit à seed donné).
\end{description}

\section{Invariants (vérifiés par tests, certains par assertions runtime)}

\begin{enumerate}[nosep]
\item[I1.] $w_i \ge 0$ pour toute entité vivante, à toute phase (clamp à 0
  autorisé uniquement pour absorber l'erreur flottante $|w| <$ 1e-12).
\item[I2.] Tout contrat actif a $q \ge$ \code{q\_min} $> 0$ et lie deux
  entités \emph{vivantes} distinctes ; aucun contrat orphelin (index
  inverses cohérents avec \code{loans}).
\item[I3.] \code{claims[i]} $= \sum_{\ell=i} q$ et \code{debts[i]}
  $= \sum_{b=i} q$ exactement (reconstruction périodique comparée aux
  agrégats incrémentaux, tolérance 1e-6 relative).
\item[I4.] Création d'un prêt : $w_\ell + w_b$ inchangé (conservation du
  principal).
\item[I5.] Service d'intérêts sans défaut : $\sum_i w_i$ inchangé.
\item[I6.] Contrats et $d_0$ jamais dépréciés ; seule la phase 5 multiplie
  $w$ par $(1-\delta)$.
\item[I7.] Bilan de faillite : la somme des pertes sèches des prêteurs
  $+$ parts détruites $=$ dettes contractuelles annulées $-$ résiduel
  distribué (loggé par cascade, vérifié en test).
\item[I8.] Fin de pas : aucune vivante avec $\NW < -$\code{tol\_nw}.
\item[I9.] Aucune borne exogène sur $N$ : seuls garde-fous d'\emph{arrêt}
  (extinction $N=0$, explosion $N > 3\cdot 10^4$) qui terminent la
  simulation avec un statut explicite, sans modifier la dynamique.
\item[I10.] Deux runs de même config et seed produisent des séries et
  snapshots identiques bit à bit.
\end{enumerate}

\section{Risques numériques identifiés}
(i)~Prolifération de contrats par scission C5 : suivi du nombre de contrats
par pas, seuil d'alerte ; les parts $<$ \code{q\_min} sont abandonnées.
(ii)~$\sqrt{w}$ pour $w$ minuscule : sans danger ; $r^* = \alpha/(2\sqrt w)$
pour $w \to 0$ : gardé par C4. (iii)~Accumulation d'erreur flottante dans
\code{claims}/\code{debts} incrémentaux : reconstruction de contrôle en fin
de run (et dans les tests), pas à chaque pas. (iv)~Cascades longues :
plafond d'itérations $=$ nombre de vivantes $+ 1$ (une faillite au moins par
itération sinon point fixe) — dépassement $=$ bug, exception.

\section{Conventions de temps et d'unités}
$\alpha = 1$ fixe l'unité de flux ; $\delta$ est l'inverse de l'unité de
temps (un pas $=$ un « pas de marché » : chaque vivante échantillonnée en
moyenne une fois par le marché) ; $\lambda$ fixe l'échelle de $N^*$.
Transitoire exclu des mesures : $t < 500$. Fenêtres de mesure :
$[500,1000)$, $[1000,1500)$, $[1500,2000)$, instantanés poolés tous les 50
pas \emph{au sein d'une fenêtre, d'un seed} ; jamais de pooling inter-seed.

\section{Sorties}
\paragraph{Séries par pas (CSV \code{series.csv}).} \code{t, births, deaths,
pop, w\_tot, nw\_tot, n\_loans, loan\_volume\_new, interest\_paid,
defaults, cascade\_iters, top\_decile\_deaths}.
\paragraph{Snapshots (NPZ \code{snap\_t<t>.npz}).} Vecteurs alignés
\code{id, w, claims, debts, nw, income, income\_net, age} des vivantes.
Cadence : tous les 50 pas à partir de $t=500$, plus $t \in \{100, 250\}$
pour le transitoire.
\paragraph{Config et résumé.} \code{config.json} (dataclass + seed + version
du code) ; \code{summary.json} (statut de fin, durée, compteurs).
\paragraph{Figures (PNG 150 dpi, \code{figures/}).} Générées par
\code{plots.py} à partir des fichiers ci-dessus (jamais depuis l'état en
mémoire) : séries temporelles ; histogrammes/CCDF log-log de $w$, $\NW$,
revenu ; QQ-plot exponentiel du corps ; âge vs $\log \NW$ ; distribution des
tailles de cascade.

\section{Protocole de validation (résumé opérationnel)}
\begin{enumerate}[nosep]
\item \textbf{Baseline} : config du rapport ($\delta=0{,}05$, $k=6$,
  $\sigma=0{,}25$, $\lambda=10$, $w_0=30$, $d_0=28$, $N(0)=0$, $T=2000$),
  3 seeds.
\item \textbf{Corps} (sur $\NW$, $y$, $w$ ; par fenêtre, par seed) : MLE
  exponentielle vs gamma, log-normale, Fisk sur les 95\,\% inférieurs
  (AIC/BIC) ; QQ-plot ; med/mean ; test mécanistique $T \approx B_0/A_0$ sur
  les incréments de $\NW$ des entités du corps (fenêtre de $\NW$ fixe).
\item \textbf{Queue} : CSN ($x_{\min}$ par KS, MLE, bootstrap $\ge 200$
  rééchantillonnages), LR Vuong contre log-normale, par fenêtre et par
  seed ; exposants comparés à $x_{\min}$ commun (médiane des $x_{\min}$).
\item \textbf{Anti-cohorte} : corr(âge, $\log \NW$) ; âges par décile de
  $\NW$ ; CSN à âge contrôlé ($[50,150)$, $[150,400)$) \emph{sur $\NW$} ;
  renouvellement du top décile entre fenêtres (fraction née récemment,
  survie, mortalité du top).
\item \textbf{SOC} : distribution des faillites par pas et par $k \in
  \{2,3,4,6\}$ ; concentration du crédit (part du top décile des créances)
  vs taille de cascade ; réponse de l'exposant à
  $\sigma \in \{0{,}15, 0{,}25, 0{,}35\}$.
\item \textbf{Grille de robustesse} : $\sigma \in \{0{,}15,0{,}25,0{,}35\}
  \times \varepsilon_0/w_0 \in \{0{,}05,0{,}25,0{,}5\} \times k \in \{3,6\}$
  ($w_0 = 30$ fixe, $d_0$ ajusté), 1 seed/case.
\item \textbf{Ablations} : \code{fail\_lender\_loans=cancel} (réplication du
  prototype E7c), \code{fail\_residual=destroy} (contrôle anti-condensation),
  \code{credit=False} (modèle nul GBM+extraction+$d_0$).
\end{enumerate}

\section{Tests unitaires (liste normative)}
Création d'entité (dotation, $d_0$, $\NW = \varepsilon_0$) ; calcul de
$\NW$ avec contrats ; création de contrat (conservation I4, index I2/I3) ;
service complet (I5) ; service partiel (transfert du disponible, défaut,
$w_b = 0$) ; dépréciation (I6) ; règle de marché (cas $q_{\mathrm{dem}}$
limitant, $q_{\mathrm{off}}$ limitant, $w_\ell \le w_b$, $N < k$, $N \in
\{0,1\}$) ; faillite simple (annulation, perte sèche, transfert C5,
résiduel C6, $d_0$ éteinte) ; cascade à deux niveaux construite à la main ;
part au borrower-créancier (annulation C5) ; créancier mort (C7) ;
invariants I1--I3 après un run court ; reproductibilité I10 ; cas limites
$w = 0$, $q \le$ \code{q\_min} ; garde-fous I9.

\end{document}
